\documentclass[11pt]{article}

\usepackage[utf8]{inputenc}
\usepackage[T1]{fontenc}
\usepackage[spanish,es-noquoting,es-tabla]{babel}
\usepackage[letterpaper,margin=2.2cm]{geometry}
\usepackage{newtxtext,newtxmath}
\usepackage{booktabs}
\usepackage{enumitem}
\usepackage{xcolor}
\usepackage{mdframed}
\usepackage{titlesec}
\usepackage[hidelinks]{hyperref}
\usepackage{siunitx}
\sisetup{output-decimal-marker={,},group-separator={{.}}}

\definecolor{azul}{HTML}{0B3D91}
\definecolor{rojo}{HTML}{C0392B}
\definecolor{gris}{HTML}{555555}
\definecolor{fondo}{HTML}{F4F6F9}

\titleformat{\section}{\large\bfseries\color{azul}}{\thesection.}{0.6em}{}
\titleformat{\subsection}{\normalsize\bfseries}{\thesubsection}{0.6em}{}
\setlist{nosep, leftmargin=1.2em}

% Bloque para la respuesta preparada a cada pregunta.
\newmdenv[backgroundcolor=fondo, linewidth=0pt, skipabove=4pt, skipbelow=8pt,
          innertopmargin=6pt, innerbottommargin=6pt]{resp}

% Pregunta del jurado.
\newcommand{\pregunta}[1]{\medskip\noindent\textbf{\color{rojo}P.}\; \textit{#1}\par}
\newcommand{\clave}[1]{\noindent\textbf{#1}}

\title{\vspace{-1.5cm}\textbf{Guion de defensa}\\[2mm]
\large Conteo de rocas visibles basado en las máscaras de AI4Mars}
\author{Juan Pablo Delgado Castro}
\date{}

\begin{document}
\maketitle
\thispagestyle{empty}

\noindent\rule{\textwidth}{0.4pt}
\begin{center}
\textit{Documento de trabajo. No forma parte del documento de tesis.}
\end{center}
\noindent\rule{\textwidth}{0.4pt}

\vspace{4mm}

\section{El hilo: qué contar y en qué orden}

La defensa es de la \textbf{pregunta de investigación}. Todo lo demás —la aplicación, el
sistema de avisos— es soporte, y solo aparece al final y en un minuto.

\subsection*{1. El problema (2 min)}

Las rocas condicionan lo que una misión puede hacer: dañan ruedas, bloquean rutas y entran
en la selección de sitio de aterrizaje. Existe una tradición consolidada de medirlas, pero
es \textbf{manual y sobre pocas escenas}.

AI4Mars aporta máscaras de segmentación por píxel para decenas de miles de imágenes,
hechas por voluntarios. Y ese material se ha usado casi exclusivamente para \textbf{entrenar
redes}: la máscara es la respuesta correcta contra la que se mide un modelo.

\clave{La observación de partida:} esa máscara no es solo la respuesta de un problema de
aprendizaje. Es, en sí misma, un mapa cuantitativo del terreno. Si un píxel dice «aquí hay
roca», contar píxeles mide cuánta roca hay.

\subsection*{2. La pregunta (30 s)}

Cómo cuantificar, a partir de esas máscaras, la cobertura de roca visible y el número
aproximado de rocas por imagen, con parámetros explícitos y resultados reproducibles.

\clave{Decir en voz alta que la pregunta tiene dos mitades de dificultad muy distinta.}
Medir cobertura es contar píxeles con un denominador bien elegido. Contar rocas exige
separar instancias sobre una máscara que no distingue instancias.

\subsection*{3. El método (3 min)}

Una frase por etapa, sin detenerse: máscara binaria, limpieza morfológica, componentes
conectadas, transformada de distancia, división de aguas guiada por máximos de prominencia,
filtros de área y forma.

\clave{Lo que hay que subrayar:} todos los umbrales son números explícitos, documentados y
discutibles. Corre en un portátil, sin entrenar nada. Esa transparencia es el argumento del
método, no su sofisticación.

\subsection*{4. Los resultados (5 min)}

\begin{enumerate}
  \item \textbf{La cobertura funciona.} Calculada para las \num{16064} escenas; definida en
    el \SI{67.3}{\percent}. Se comprobó que los valores altos no son un artefacto de la
    fórmula ($r = -0{,}02$ frente a la fracción etiquetada).
  \item \textbf{Los dos indicadores son independientes} ($r = 0{,}02$). No es trivial:
    significa que miden cosas distintas y ninguno sustituye al otro.
  \item \textbf{El conteo tiene un techo.} Validado contra conteo humano en \num{60}
    escenas: en ninguna el procedimiento supera las nueve rocas.
\end{enumerate}

\subsection*{5. Los tres hallazgos que sostienen el trabajo (5 min)}

Esta es la parte que hay que contar bien. Son resultados, no limitaciones.

\begin{enumerate}
  \item \textbf{La anotación colaborativa sobreestima la roca.} La cobertura mediana pasa
    de \SI{96.8}{\percent} a \SI{46.1}{\percent} con máscaras de experto. \emph{Esto afecta
    a todo el que use AI4Mars, no solo a mí.}
  \item \textbf{El techo del conteo es la taxonomía, no las imágenes.} Las máscaras de
    experto tienen \emph{menos} roca grande, no más. No se arregla con mejores datos.
  \item \textbf{La anotación no es exhaustiva.} Marca algunas rocas, no todas: mediana del
    \SI{9.1}{\percent} de la roca etiquetada. El indicador mide bloques dentro de lo que el
    anotador delimitó, no rocas de la escena.
\end{enumerate}

\subsection*{6. Cierre (1 min)}

El trabajo entrega un procedimiento auditable, un conjunto de resultados reutilizable y un
diagnóstico cuantitativo de los límites del material. Y deja construida una capa aplicada
—avisos de terreno y una aplicación de consulta— que muestra que los indicadores son
utilizables. \textbf{Un minuto, no más.}

\newpage
\section{Cifras que hay que tener en la cabeza}

\begin{table}[h]
\centering\small
\begin{tabular}{lr}
\toprule
\multicolumn{2}{l}{\textbf{Conjunto}}\\
Escenas analizadas & \num{16064}\\
Con roca visible & \SI{67.3}{\percent}\\
Con roca grande & \SI{13.9}{\percent}\\
Fracción de escena etiquetada (mediana) & 0,58\\
\midrule
\multicolumn{2}{l}{\textbf{Cobertura (E1)}}\\
Mediana sobre píxeles etiquetados & \SI{96.8}{\percent}\\
Mediana sobre la imagen completa & \SI{42.0}{\percent}\\
Correlación con la fracción etiquetada & $-0{,}02$\\
\midrule
\multicolumn{2}{l}{\textbf{Conteo (E2)}}\\
Escenas aptas / rocas contadas & \num{2193} / \num{4204}\\
Mediana por imagen / máximo & 1 / 20\\
Correlación con la cobertura & $0{,}02$\\
\midrule
\multicolumn{2}{l}{\textbf{Sesgo de la anotación}}\\
Cobertura mediana: colaborativa $\to$ experto & \SI{96.8}{\percent} $\to$ \SI{46.1}{\percent}\\
Escenas con cobertura total & \SI{42}{\percent} $\to$ \SI{8}{\percent}\\
Píxeles de suelo y arena & \SI{50}{\percent} $\to$ \SI{69}{\percent}\\
\midrule
\multicolumn{2}{l}{\textbf{Techo del conteo}}\\
Roca etiquetada que es \textit{bedrock} & \SI{97.5}{\percent}\\
Escenas sin ningún píxel de roca grande & \num{13838} (\SI{86.1}{\percent})\\
Roca grande en experto (niveles 1/2/3) & \SI{16.5}{\percent} / \SI{6.2}{\percent} / \SI{1.6}{\percent}\\
\midrule
\multicolumn{2}{l}{\textbf{Validación humana ($n = 60$)}}\\
Acuerdo exacto & \SI{50}{\percent}\\
Kappa / kappa ponderado & 0,13 / 0,11\\
Por debajo / por encima del observador & 28 / 2\\
Conteo máximo del procedimiento & 9\\
Roca grande sobre la roca etiquetada (mediana) & \SI{9.1}{\percent}\\
\midrule
\multicolumn{2}{l}{\textbf{Comparación con aprendizaje automático}}\\
DeepLabV3, IoU medio (colaborativas / experto) & 0,940 / 0,843\\
Correlación de la cobertura (colab. / experto) & 0,950 / 0,927\\
Modelo general (FastSAM), acuerdo por bandas & \SI{52}{\percent}\\
\bottomrule
\end{tabular}
\end{table}

\section{Qué NO afirmar}

Lo más peligroso en una defensa no es no saber algo: es afirmar de más y que te lo
desmonten. Estas cinco frases están prohibidas.

\begin{itemize}\itemsep3pt
  \item \textbf{«Mi sistema mide la abundancia de roca en Marte.»} Mide la abundancia
    \emph{en las máscaras de AI4Mars}, que sobreestiman. Decirlo siempre en relativo.
  \item \textbf{«Cuento las rocas de la imagen.»} Cuento bloques dentro de la región que el
    anotador marcó como roca grande, que es una fracción pequeña de la roca visible.
  \item \textbf{«Esto sirve para conducir el rover.»} El procedimiento necesita una máscara
    que no existe en el momento de la bajada de datos. Lo que sí puede correr sin humanos
    es el segmentador entrenado, y eso es una línea, no un resultado operativo.
  \item \textbf{«El método clásico es mejor que aprendizaje profundo.»} No se compararon en
    la misma tarea. El clásico es auditable y barato; el entrenado estima cobertura sin
    máscara. Son complementarios.
  \item \textbf{«La validación demuestra que el conteo es correcto.»} Un observador,
    \num{60} escenas, kappa 0,13. Demuestra lo contrario: que tiene un techo.
\end{itemize}

\section{Banco de preguntas difíciles}

\subsection*{Sobre los resultados}

\pregunta{Una cobertura mediana del 96,8\,\% implicaría que Marte es casi todo roca. ¿No
es eso implausible?}
\begin{resp}
Lo es, y por eso lo investigué. Ese valor es relativo a los \emph{píxeles etiquetados}, no
a la escena. La comprobé contra las \num{322} máscaras de experto que el propio conjunto
incluye, con el mismo procedimiento y sin tocar un solo parámetro: la mediana cae al
\SI{46.1}{\percent}. El sesgo no está en mi cálculo —aplicado a las máscaras de experto da
valores plausibles— sino en el dato de entrada. Por eso reporto también la cobertura sobre
la imagen completa, que es una cota inferior conservadora.
\end{resp}

\pregunta{¿Y cómo sabe que las coberturas altas no son simplemente escenas con pocos
píxeles etiquetados?}
\begin{resp}
Lo contrasté explícitamente: la correlación entre la cobertura y la fracción de escena
etiquetada es de $-0{,}02$, prácticamente nula. Las coberturas altas corresponden a escenas
genuinamente dominadas por lecho rocoso, no a un artefacto del denominador. Ese control
está en el documento como figura.
\end{resp}

\pregunta{Su conteo tiene un kappa de 0,13 con el juicio humano. ¿No es eso un fracaso?}
\begin{resp}
Es un resultado, y es el que permite decir algo que antes no se sabía. Lo importante no es
el valor sino su estructura: en ninguna de las \num{60} escenas el procedimiento supera las
nueve rocas, mientras el observador identificó diez o más en doce. Eso no es un umbral mal
puesto, es un techo. Y localicé la causa: la anotación marca como roca grande solo una
mediana del \SI{9.1}{\percent} de la roca etiquetada. Ni un procedimiento perfecto contaría
lo que la anotación no delimita.

Añadiría que el conteo era el segundo objetivo; el principal, la cobertura, sí funciona y
está definida para dos tercios del conjunto.
\end{resp}

\pregunta{¿Por qué debemos creer que la anotación está sesgada y no su método?}
\begin{resp}
Porque no lo comprobé contra un criterio propio, sino contra las máscaras de experto que el
propio conjunto de datos distribuye para validación, ejecutando el mismo código sin
modificar ningún parámetro. Si el sesgo fuera mío, aparecería también sobre las máscaras de
experto, y no aparece.
\end{resp}

\subsection*{Sobre el método}

\pregunta{¿Por qué no entrenó una red, que es el estado del arte?}
\begin{resp}
Por tres razones, y de hecho sí entrené una como comparación. Primera, transparencia: cada
umbral de mi procedimiento es un número que se puede discutir; una predicción de red no.
Segunda, costo: corre en un portátil sin conjunto de entrenamiento. Tercera, y la que
resultó más productiva: el procedimiento clásico sirve de instrumento para \emph{auditar el
dato}, que es de donde salió el hallazgo del sesgo.

Como contraste entrené un DeepLabV3 y probé un modelo fundacional. El entrenado alcanza IoU
medio de 0,940 contra las máscaras colaborativas, pero eso mide parecido con la anotación
de la que aprendió, no acierto.
\end{resp}

\pregunta{¿Cómo eligió los umbrales? ¿No son arbitrarios?}
\begin{resp}
Se calibraron por inspección visual sobre escenas de terrenos distintos, comparando el corte
con la imagen original, y quedan documentados con su efecto. Además comprobé qué se pierde
por ellos: en \num{453} escenas los filtros dejan el conteo en cero, las revisé caso por
caso y \emph{aciertan} —son motas de doce píxeles o bandas del horizonte—. Los umbrales de
los avisos se fijaron sobre percentiles de la distribución observada, no a ojo.
\end{resp}

\pregunta{Dice que las imágenes no son el problema, pero su procedimiento ni siquiera las
abre. ¿Cómo lo sabe?}
\begin{resp}
Es exactamente la objeción correcta, y la respondí con tres medidas. Primera: las máscaras
de experto tienen \emph{menos} roca grande, no más, así que no es un problema de calidad de
anotación ni de imagen. Segunda: probé un enfoque híbrido que sí usa la imagen dentro de la
región anotada; rompe el techo, pero la mejora en acuerdo no es estadísticamente
demostrable con \num{60} escenas. Tercera: el segmentador entrenado lee las imágenes
directamente y alcanza IoU de 0,843 contra etiquetas de experto, lo que prueba que la
señal está en la imagen.
\end{resp}

\subsection*{Sobre la validación}

\pregunta{Su validación tiene un solo observador. ¿Qué valor tiene?}
\begin{resp}
Es una limitación real y está declarada en el documento. No permite estimar la variabilidad
entre personas, que en esta tarea es previsiblemente alta, y por eso lo reporto como
comprobación cualitativa y no como medición con intervalos de confianza.

Dicho eso, el hallazgo principal no depende de la finura del juicio: que el procedimiento
no supere nunca las nueve rocas mientras el observador identifica veinticinco es una
diferencia de orden de magnitud, no de apreciación.
\end{resp}

\pregunta{¿No pudo diseñar la validación para que le diera el resultado que quería?}
\begin{resp}
Tomé precauciones explícitas contra eso: identificadores neutros, orden barajado, y el
resultado del algoritmo no aparece en ningún momento del material que se consulta al
responder. La muestra se estratificó por área de la región anotada, deliberadamente
\emph{no} por la banda del algoritmo, para no condicionarla al método evaluado.

Hay además una evidencia de que no lo diseñé a favor: una primera ronda de \num{24}
escenas resultó defectuosa —su banda superior premiaba a los métodos que sobreestiman— y la
rehíce entera. El resultado empeoró, y aun así es el que reporto.
\end{resp}

\subsection*{Sobre el alcance y el aporte}

\pregunta{¿No es esto simplemente estadística descriptiva sobre etiquetas ajenas?}
\begin{resp}
El cálculo de la cobertura sí es, en esencia, un conteo bien definido, y no pretendo otra
cosa. El aporte no está ahí. Está en haber medido que esas etiquetas —que toda una línea de
investigación usa como referencia para evaluar modelos— sobreestiman la roca de forma
sistemática, y en cuantificar cuánto. Cualquier trabajo que reporte índices de solapamiento
contra AI4Mars está midiendo el acuerdo con una anotación sesgada, y eso no estaba
documentado.
\end{resp}

\pregunta{¿Para qué sirve todo esto en la práctica?}
\begin{resp}
Para tres cosas, en orden decreciente de solidez. La cobertura permite comparar escenas
dentro del conjunto y caracterizar tramos de recorrido, y eso funciona. El conteo distingue
de forma aproximada escenas con y sin bloques discretos, con un sesgo conocido a la baja. Y
el diagnóstico del sesgo sirve a cualquiera que entrene sobre este conjunto.

Construí además una capa de avisos de terreno y una aplicación de consulta, pero las
presento como demostración de que los indicadores son utilizables, no como un producto
validado: heredan las limitaciones de los indicadores de los que derivan.
\end{resp}

\pregunta{¿Qué haría diferente si empezara hoy?}
\begin{resp}
Tres cosas. Habría verificado la codificación de clases desde el primer día: la que asumí
en el anteproyecto era incorrecta y de no haberla comprobado habría invalidado todo en
silencio. Habría diseñado la validación humana con más observadores y desde el principio con
la escala correcta. Y habría mirado antes la taxonomía geológica extendida del conjunto, que
distingue roca suelta y es la vía natural para sortear el techo del conteo.
\end{resp}

\pregunta{¿Y si el jurado insiste en que el conteo no funciona?}
\begin{resp}
Concederlo sin defenderlo. «Tiene razón: el conteo no reproduce el juicio humano, y lo
reporto así. Lo que aporta el trabajo no es un contador de rocas, sino la medición de por
qué no se puede construir con este dato, que es información que no estaba disponible.»
\end{resp}

\section{Si algo falla}

\begin{itemize}\itemsep3pt
  \item \textbf{Si no sabe una respuesta:} «No lo medí, y no quiero improvisar una cifra.
    Lo que sí puedo decir es\dots». Un jurado castiga mucho más inventar que reconocer.
  \item \textbf{Si le cuestionan una cifra:} está toda en el documento con su figura o su
    tabla. Remitir a ella.
  \item \textbf{Si se enredan en lo técnico:} volver a la pregunta de investigación. La
    defensa es de ella, no de la implementación.
\end{itemize}

\end{document}
